\textbf{1. Kiến thức cần nhớ}
\begin{itemize}
\item \textbf{Phân tích số liệu}: Từ bảng số liệu hoặc biểu đồ, ta có thể:
\begin{itemize}
\item Tìm giá trị lớn nhất (cao nhất), giá trị nhỏ nhất (thấp nhất).
\item Tính giá trị trung bình, tỉ lệ tăng/giảm giữa các thời điểm.
\item Nhận xét xu hướng phát triển chung (tăng trưởng, suy thoái, ổn định).
\end{itemize}
\item \textbf{Phát hiện số liệu sai lệch}:
\begin{itemize}
\item Trục tung không bắt đầu từ số /bin/zsh$ hoặc chia tỉ lệ không đều làm phóng đại hoặc thu nhỏ sai lệch thực tế.
\item Hình vẽ biểu diễn diện tích/thể tích không tỉ lệ với chiều cao số liệu.
\item Tổng các hình quạt trong biểu đồ quạt tròn khác \%$.
\end{itemize}
\end{itemize}

\textbf{2. Phương pháp giải các dạng}

%% DANG 8D5?20-1
\textbf{Dạng 1. Đọc và phân tích số liệu từ biểu đồ}
\begin{enumerate}
\item Quan sát kĩ trục nằm ngang (đối tượng/thời gian) và trục thẳng đứng (độ lớn, đơn vị).
\item Dóng các điểm/cột sang trục tung để đọc số liệu chính xác.
\item Thực hiện các phép so sánh hơn kém, tính chênh lệch: $\text{Chênh lệch} = \text{Giá trị sau} - \text{Giá trị trước}$.
\item Tính phần trăm tăng/giảm:
80830\% \text{ tăng/giảm} = \dfrac{|\text{Giá trị sau} - \text{Giá trị trước}|}{\text{Giá trị trước}} \times 100\%80830
\end{enumerate}

%% DANG 8D5?20-2
\textbf{Dạng 2. Phát hiện biểu đồ, số liệu sai lệch}
\begin{enumerate}
\item Kiểm tra gốc tọa độ trục tung có phải số /bin/zsh$ không.
\item Kiểm tra các khoảng chia trên trục tọa độ có đều nhau không.
\item Kiểm tra tổng tỉ lệ phần trăm trên biểu đồ quạt tròn có đúng bằng \%$ không.
\item Đưa ra nhận xét phê phán tính trung thực của biểu đồ khi có sự bóp méo trực quan.
\end{enumerate}

%% DANG 8D5?20-3
\textbf{Dạng 3. Rút ra kết luận, giải quyết vấn đề thực tế}
\begin{enumerate}
\item Từ xu hướng biến thiên của số liệu, đưa ra dự báo hợp lí cho tương lai gần.
\item Đề xuất giải pháp thực tế dựa trên số liệu thống kê (sản xuất loại sản phẩm nào nhiều hơn, lên kế hoạch tài chính, phòng chống dịch bệnh\dots).
\end{enumerate}
